% ============================================================
%  BAB 4 – PENUTUP
% ============================================================


\chapter*{PENUTUP}
\refstepcounter{chapter}
\addcontentsline{toc}{chapter}{PENUTUP}
\pdfbookmark[0]{PENUTUP}{penutup}


\section{Kesimpulan}

Berdasarkan hasil penelitian dan pembahasan yang telah dilakukan, diperoleh kesimpulan
sebagai berikut.
\begin{enumerate}
  \item Prototipe aplikasi Dalsilens berhasil dikembangkan sebagai alat bantu visual
        berbasis Android secara \textit{real-time}.
        Aplikasi mampu mengakuisisi citra dari kamera menggunakan CameraX,
        memproses setiap \textit{frame} secara langsung, dan menampilkan hasil
        transformasi warna tanpa jeda yang berarti bagi pengguna.
        Pengujian fungsionalitas menggunakan metode \textit{Black Box Testing}
        terhadap 15 skenario menunjukkan seluruh fitur berjalan sesuai ekspektasi
        dengan hasil \textbf{Berhasil} pada semua skenario.

  \item Algoritma \textit{LMS Daltonization} berhasil diimplementasikan untuk
        membantu penderita dikromasi membedakan objek secara visual.
        Matriks transformasi warna dihitung melalui tiga tahap konversi ruang warna:
        RGB $\rightarrow$ LMS $\rightarrow$ simulasi defisiensi $\rightarrow$ RGB,
        kemudian diaplikasikan ke \textit{frame} kamera menggunakan
        \textit{ColorMatrix} Android yang diakselerasi GPU.
        Fitur ini mendukung tiga jenis dikromasi, yaitu \textit{Protanopia},
        \textit{Deuteranopia}, dan \textit{Tritanopia}, dengan parameter
        \textit{severity} yang dapat disesuaikan.

  \item Fitur deteksi nama warna berbasis ruang warna HSV berhasil
        diimplementasikan untuk membantu pengguna membedakan objek secara tekstual.
        Klasifikasi dilakukan menggunakan pendekatan \textit{rule-based thresholding}
        yang membagi ruang warna menjadi 10 kategori: tujuh warna kromatik
        (merah, oranye, kuning, hijau, cyan, biru, ungu) dan tiga warna akromatik
        (putih, abu-abu, hitam).
        Hasil pengujian akurasi pada kondisi pencahayaan normal mencapai
        \textit{[X]}\% dan pada kondisi pencahayaan redup mencapai \textit{[Y]}\%.
\end{enumerate}

\section{Saran}

Berdasarkan penelitian yang telah dilakukan, saran untuk pengembangan selanjutnya
adalah sebagai berikut.
\begin{enumerate}
  \item Fitur klasifikasi nama warna saat ini menggunakan pendekatan
        \textit{rule-based thresholding} dengan nilai batas yang ditentukan secara
        empiris.
        Pengembangan selanjutnya dapat menggunakan pendekatan \textit{machine
        learning} berbasis data untuk meningkatkan akurasi deteksi, terutama pada
        kondisi pencahayaan yang berubah-ubah.

  \item Aplikasi saat ini hanya mendukung tiga jenis dikromasi (buta warna total
        pada satu kanal).
        Pengembangan berikutnya dapat mencakup tipe anomalous trichromacy seperti
        \textit{Protanomaly} dan \textit{Deuteranomaly} yang jumlah penderitanya
        lebih banyak di populasi umum.

  \item Antarmuka pengguna dapat dikembangkan dengan menambahkan mode
        \textit{accessibility} yang lebih lengkap, seperti umpan balik suara
        (\textit{text-to-speech}) untuk nama warna yang terdeteksi, sehingga
        aplikasi dapat digunakan oleh pengguna dengan keterbatasan visual ganda.
\end{enumerate}
